\begin{table}[htbp]
\centering
\caption{Pemex y CFE: por dónde entra y en qué se va (mdp)}
\label{tab:C9_2}
\small
\begin{tabular}{lSSSSS}
\toprule
{Concepto} & {2024} & {2025} & {Dif. nominal} & {Dif. real} & {Var. real \%} \\
\midrule
Pemex: gasto total (bruto) & \num{624805.6} & \num{612145.9} & \num{-12659.7} & \num{-39526.3} & \num{-6.1} \\
Pemex: pensiones (tipo de gasto 4) & \num{74120.5} & \num{83827.6} & \num{9707.1} & \num{6519.9} & \num{8.4} \\
Pemex: obra publica (capitulo 6000) & \num{237574.8} & \num{199216.9} & \num{-38357.9} & \num{-48573.6} & \num{-19.6} \\
Pemex: servicios personales (capitulo 1000) & \num{105084.2} & \num{114074.0} & \num{8989.8} & \num{4471.2} & \num{4.1} \\
Pemex: deuda publica (capitulo 9000) & \num{143341.2} & \num{147890.7} & \num{4549.5} & \num{-1614.2} & \num{-1.1} \\
CFE: gasto total (bruto) & \num{528532.7} & \num{583951.1} & \num{55418.4} & \num{32691.5} & \num{5.9} \\
CFE: pensiones (tipo de gasto 4) & \num{55961.3} & \num{64337.9} & \num{8376.6} & \num{5970.3} & \num{10.2} \\
CFE: obra publica (capitulo 6000) & \num{32405.2} & \num{39064.5} & \num{6659.3} & \num{5265.9} & \num{15.6} \\
CFE: servicios personales (capitulo 1000) & \num{74169.5} & \num{80447.5} & \num{6278.0} & \num{3088.7} & \num{4.0} \\
CFE: deuda publica (capitulo 9000) & \num{35152.0} & \num{38479.5} & \num{3327.5} & \num{1816.0} & \num{5.0} \\
Ramo 18 Energia (aportacion del Gobierno Federal) & \num{167736.2} & \num{138307.4} & \num{-29428.8} & \num{-36641.5} & \num{-20.9} \\
\bottomrule
\end{tabular}
\par\vspace{2pt}\footnotesize Fuente: elaboración propia del ITED con los analíticos del PEF aprobado 2024 y 2025. Las diferencias en millones de pesos se reportan dos veces, nominal y real; la real está en pesos de 2025 con el deflactor de 4.3\% que declara el CGPE.
\end{table}
